\section{実験プロトコル}
\label{sec:experiments}

本稿では、実装上の主張と、実験結果に基づく採用可否の主張とを分離する。
\cref{sec:model,sec:consistency}のモデルと損失によって本手法は定義されるが、
これらを実運用のデフォルトとして選択するには、以下に示す固定条件下での
複数seed比較が必要である。この分離により、整合性残差が小さいことを、
いずれかの分岐の精度が向上した証拠として報告することを防ぐ。

\subsection{データ分割と統制変数}

主要な学習データ源には、\cref{sec:render_supervision}で説明した、対象コート単位の
singleton合成データセットを用いる。隣接するレンダリング画像が異なる分割にまたがらないよう、
軌道グループ全体を80/10/10に分割する。検証用およびテスト用manifestには、視点の方位角、
仰角、画角、coverage mode、対象コート、renderer-support flag、および軌道識別子を
記録しなければならない。このメタデータにより、難しい部分視野を平均値の中に埋没させる
ことなく層別解析できる。

人手でアノテーションしたholdout実画像集合を用いて、\KP{}、ライン、および
セグメンテーションの転移性能を評価する。独立したキャリブレーション正解情報が
与えられない限り、この集合では並進、回転、焦点距離を評価できない。また、評価対象の
検出器から擬似姿勢ラベルを作ることはしない。したがって、合成データ上の姿勢指標と
実画像上の密予測指標は、単一のスコアに統合せず、別々の列に報告する。

同一のablation系列では、合成データのmanifest、\DINO{} checkpointとtrain mode、
姿勢を損なわないaugmentation、optimizer、batch構成、直接損失の重み、画像scale
（短辺約256 pixels）、および15-epochのbudgetを固定する。最終比較にはそれぞれ少なくとも
3つのseedを用いる。temperature、Huber threshold、warmup、またはauxiliary weightの
hyperparameter探索は、短いseed-0 pilotに限定し、選択した値は複数seedでの実行前に固定する。

\subsection{順序を定めたアーキテクチャと目的関数のアブレーション}

幾何学的な補助損失を比較する前に、アーキテクチャを選択する。まずdecoderを固定して
task encoderのdepthをsweepし、次に選択したencoderを用いてdecoder familyとcapacityを
sweepする。その後に初めて、アーキテクチャを固定してsupervisionを比較する。この順序により、
capacityの変化をconsistency objectiveの効果と誤って解釈することを避ける。

\begin{table}[t]
  \centering
  \caption{Ordered ablation plan.  Every final row is repeated with at least
  three seeds after pilot hyperparameters are frozen.}
  \label{tab:ablation_protocol}
  \small
  \begin{tabularx}{\linewidth}{@{}p{0.14\linewidth}p{0.24\linewidth}p{0.25\linewidth}X@{}}
    \toprule
    Stage & Frozen & Varied & Required output \\
    \midrule
    Encoder & Backbone, decoder, direct losses & Task-Transformer depth/width &
      KP/pose accuracy, variance, parameters, time \\
    Decoder & Selected encoder, supervision & linear / progressive / DPT,
      width, taps & Dense accuracy--capacity--speed Pareto \\
    Supervision & Selected architecture, all direct losses & auxiliary off,
      both-grad, two stop-gradient directions & Every GT metric, consistency,
      invalid depth, gradients, memory/time \\
    Robustness & Selected supervision & coverage/pose strata and local-context
      diagnostic & Mean and worst-stratum synthetic/real transfer \\
    \bottomrule
  \end{tabularx}
\end{table}



中核となるsupervisionの比較は以下のとおりである。
\begin{enumerate}
  \item \textbf{direct-all}: \KP{}、ライン、セグメンテーション、並進、回転、焦点距離に
    対する直接損失。
  \item \textbf{joint-both}: direct-allに\cref{eq:cheirality}を加え、両方の分岐へ
    勾配を伝播する。
  \item \textbf{joint-stopgrad-pose}: 姿勢によるprojectionの信号を密予測分岐のみに伝える。
  \item \textbf{joint-stopgrad-dense}: 密予測の期待値の信号を姿勢分岐のみに伝える。
\end{enumerate}
4つの設定はすべて同一の直接教師信号を維持する。stop-gradient variantは、どちらの
分岐が恩恵を受けるかを診断するものであり、test-time graphの代替ではない。推論時には
補助経路を取り除くため、4つのvariantでdeployment latencyとparameter countは同一である。

\subsection{評価指標と循環性のない評価基準}

密予測キーポイントは、モデルの画素座標系におけるEuclidean distanceの平均値と中央値で評価し、
channel別およびsupervision eligibility別に層別した内訳も示す。ラインにはDice、
セグメンテーションにはmean intersection over unionを用いる。姿勢には、metres単位の
camera-centre L2 error、degrees単位の\(\SOthree\) geodesic error、focal relative error、
およびabsolute log-focal errorを用いる。

2つのprojection指標を混同してはならない。\emph{姿勢再投影誤差}は、予測姿勢と
予測focalのみを用いてground-truth canonical pointをprojectし、ground-truth image pointと
比較する。これは独立した姿勢指標である。\emph{分岐間整合距離}は、密予測の期待値と
予測姿勢によるprojectionを比較する。有用な診断指標ではあるが、ground-truth accuracyの
指標ではない。両分岐がともに誤っていても小さくなり得るためである。さらに、
supervision-eligible count、invalid predicted-depth rate、両分岐のfinite gradient status、
train-step time、およびpeak accelerator memoryを記録する。

\subsection{事前に宣言した採用基準}

direct-allと比較して、joint-bothは、GTキーポイント平均距離または姿勢のみの再投影誤差の
いずれかを、3-seed meanで少なくとも5\%改善しなければならない。もう一方の指標、
およびtranslation/rotation/focal errorは、相対的に5\%を超えて悪化してはならない。
line Diceとsegmentation mIoUの絶対値の低下は0.01以下でなければならない。いずれの実行でも
non-finite損失、non-finite勾配、またはall-invalid batchが生じてはならず、step timeと
peak memoryのoverheadはいずれも10\%以内でなければならない。整合距離のみの
改善は、明示的な棄却条件とする。

\subsection{ショートカットに焦点を当てた診断}

test manifestを、コートのfull、near-full、partial coverage、camera height、azimuth、
target-court distance、およびsupervision-eligible point数で層別する。平均値と併せて最悪の
stratumも報告する。また、full imageと、visible junction周辺の固定サイズの近傍のみを残して
それ以外のcontextを抑制した診断用ビューを比較する。この診断用ビューは、学習にもmodel selectionにも
使用しない。局所入力のみでも大域姿勢性能が変わらない場合、平均密予測指標が良好であっても、
意図した機構が実証されたとはみなさない。
